% chapters/11/evaluasi.tex
\bagian{Asesmen dan Instrumen Penilaian}
\label{sec:asesmen-integrasi}

\begin{asesmenbox}
\textbf{Kedudukan asesmen bab ini.} Bab ini menuntaskan \SUBCPMK{2.4} dan
menutup \CPMK{2}. Asesmennya berupa \textbf{Mini Proyek AI} yang dinilai
memakai rubrik pada Lampiran~A, dengan bobot \CPMK{2} sebesar 10 persen dari
nilai akhir mata kuliah.

\vspace{0.4em}
Mini proyek ini dikumpulkan dan dipresentasikan pada minggu kesembilan. Materi
Bab~VIII sampai Bab~XI juga termasuk bahan \textbf{Ujian Tengah Semester} pada
minggu kesepuluh.
\end{asesmenbox}

\subsection*{Instrumen 1: Rubrik Mini Proyek AI}

\begin{table}[htbp]
\centering
\caption{Rubrik penilaian Mini Proyek AI beserta indikator yang diukur}
\label{tab:rubrik-mini-proyek-bab11}
{\setlength{\tabcolsep}{4pt}\footnotesize%
\begin{tabularx}{\textwidth}{@{}L{3.4cm} C{1.2cm} C{2.1cm} Y@{}}
\toprule
\headrow \textbf{Kriteria} & \textbf{Bobot} & \textbf{Indikator} & \textbf{Deskripsi capaian} \\
\midrule
Implementasi mesin keadaan & 25\% & \ind{a} &
NPC memiliki keadaan yang jelas beserta transisi yang berjalan menurut jarak
dan pewaktu; tidak terjadi kebuntuan keadaan. \\
\addlinespace
Implementasi pencarian jalur & 25\% & \ind{a} &
Pencarian jalur menghasilkan rute yang benar pada peta berhalangan, memakai
antrean prioritas, dan heuristik yang dapat dipertanggungjawabkan. \\
\addlinespace
Penggabungan keduanya & 20\% & \ind{b} &
Pencarian jalur dipicu tepat pada keadaan mengejar, diulang secara berkala,
dan kinerjanya terjaga. \\
\addlinespace
Pengujian dan laporan & 20\% & \ind{a}, \ind{b}, \ind{c} &
Terdapat tiga skenario uji beserta angka pengukurannya, hasil pengamatan yang
objektif, dan perbaikan yang ditindaklanjuti. \\
\addlinespace
Presentasi dan kerja sama & 10\% & \ind{c} &
Demonstrasi jelas, setiap anggota memahami bagiannya, serta mampu menjawab
pertanyaan penguji. \\
\bottomrule
\end{tabularx}}
\end{table}

\subsection*{Instrumen 2: Soal Penilaian}

Jawablah dengan singkat. Kunci jawaban beserta indikator yang diukur disajikan
pada Tabel~\ref{tab:kunci-konsep-ch11}.

\begin{enumerate}
  \item Jelaskan mengapa pencarian jalur tidak dijalankan setiap bingkai, dan
        sebutkan dua syarat yang memicunya.
  \item Sebuah pencarian memeriksa 200 petak. Hitunglah beban perhitungannya
        apabila empat NPC menghitung ulang setiap bingkai pada 60 bingkai per
        detik.
  \item Sebutkan tiga segi penilaian mutu perilaku NPC beserta ukuran
        masing-masing, lalu jelaskan mengapa ketiganya bertentangan.
  \item Mengapa jalur pemain pada skenario uji harus sama ketika membandingkan
        dua rancangan?
  \item Perbaikan seperti apa yang disebut perbaikan yang baik? Sebutkan tiga
        sifatnya.
\end{enumerate}

\begin{table}[htbp]
\centering
\caption{Kunci soal penilaian Bab XI beserta indikator yang diukur}
\label{tab:kunci-konsep-ch11}
{\setlength{\tabcolsep}{4pt}\footnotesize%
\begin{tabularx}{\textwidth}{@{}C{0.8cm} C{2.1cm} Y@{}}
\toprule
\headrow \textbf{No} & \textbf{Indikator} & \textbf{Pokok jawaban} \\
\midrule
1 & \ind{b} &
Karena biaya satu kali pencarian mahal, sedangkan hasilnya dapat dipakai
berkali-kali. Pencarian dipicu ketika selang hitung ulangnya habis, dan ketika
jalurnya sudah habis dilalui. \\
\addlinespace
2 & \ind{b} &
$4 \times 60 \times 200 = 48.000$ petak setiap detik, yaitu sekitar empat puluh
delapan ribu pemeriksaan petak setiap detik. \\
\addlinespace
3 & \ind{b} &
Kelancaran diukur dengan laju bingkai, daya tanggap dengan jarak rata-rata NPC
terhadap pemain, dan beban komputasi dengan banyaknya pencarian serta petak
yang diperiksa. Ketiganya bertentangan sebab menambah daya tanggap menuntut
perhitungan yang lebih banyak, dan perhitungan yang berlebihan mengurangi
kelancaran. \\
\addlinespace
4 & \ind{a} &
Agar perbedaan hasilnya benar-benar berasal dari rancangan yang diuji, bukan
dari perbedaan jalur pemainnya. \\
\addlinespace
5 & \ind{c} &
Perbaikan yang berdasarkan pengukuran, terarah pada penyebabnya, dan disertai
cara mengukur keberhasilannya. \\
\bottomrule
\end{tabularx}}
\end{table}
